\documentclass{memoir}
\usepackage{amssymb}
\usepackage{tcolorbox}
\usepackage{amsmath}
\usepackage[italian]{babel}
\usepackage[T1]{fontenc}        % Codifica dei font migliorata
\usepackage[utf8]{inputenc}     % Supporto caratteri UTF-8 (di default in LaTeX moderni)
\usepackage{ulem}
% \usepackage{unicode-math}
%\setmathfont{Garamond Math}
%\\usepackage[varbb]{newpxmath}
%\usepackage[activate={true,nocompatibility},final,tracking=true,kerning=true]{microtype}
% microtype makes text look nicer
\usepackage[italian]{babel}     % Lingua italiana
\title{Analisi mod2}
\author{Flavio}
\usepackage{tcolorbox}
\tcbuselibrary{theorems}
\NewTcbTheorem[number within=section]{mytheo}{Teorema}%
{colback=green!5,colframe=green!35!black,fonttitle=\bfseries}{th}
\NewTcbTheorem[number within=section]{deff}{Definizione}%
{colback=green!5,colframe=green!35!black,fonttitle=\bfseries}{th}
%%\theoremstyle{definition}
\newtheorem{definition}{Definition}
\newcounter{example}[section]
\newenvironment{example}[1][]{\refstepcounter{example}\par\medskip
   \noindent \textbf{Esempio~\theexample. #1} \rmfamily}{\medskip}
\begin{document}
\maketitle

\chapter{Serie Numeriche}
\section{Introduzione}

Le Serie numeriche sono trovate in molti campi, fin dall'antichità hanno descrito problemi come il paradosso di zenone o numeri decimale con espansione infinità come \(\pi \), grazie allo sviluppo dell'analisi nel '700 è possibile farne uno studio rigorso. Ad esempio il numero 0.333:

\[0.333 = \frac{3}{10} + \frac{3}{100} + \frac{3}{1000} + \cdots + \frac{3}{10^n}
\]
\begin{deff}{Serie Numerica}{ser_num}
  Data la successione di termine generica \(a_k\), costruiamo la nuova successione \(\{s_n\}\) come:
\begin{align*}
s_0 &= a_0\\
s_1 &= a_0 + a_1\\
s_2 &= a_0 + a_1 + a_2\\
      &\vdots \\
  s_n &= a_0 + a_1 + \cdots + a_n = \sum_{k=0}^n a_k
\end{align*}
\(s_n\) è la \textbf{somma parziale n-esima} della serie.
\[
  \sum_{k=0}^{+\infty} a_n\] è la serie di termine generale \(a_n\)
\end{deff}
% TODO: fa esempio
\section{Carattere di una serie}

\subsection{Serie convergente}
Una serie \(\sum_{k=0}^{+\infty} a_n \) \textit{converge} se la successione delle somme parziali converge ad un numero \(l\), detto somma della serie.

\fbox{%
  \parbox{\textwidth}{
    \begin{example}
      Serie di Mengoli\\

      \[
        \sum_{i=0}^{+\infty} \frac{1}{n} - \frac{1}{n+1} \] \\
      \begin{align*}
        s_1 &= 1 -\frac{1}{2} \\
        s_2 &= 1-\frac{1}{2}+\frac{1}{2}-\frac{1}{3} = 1-\frac{1}{3}\\
        \vdots\\
        s_n &= 1-\frac{1}{2} + \frac{1}{2} - \frac{1}{3} = 1-1/3 + \ldots + \frac{1}{n} - \frac{1}{n-1}            \end{align*}
      Converge perchè: \(\lim_{n \to \infty} s_n = 1\)
    \end{example}
}%

}

\subsection{serie divergente}
Una serie può Divergere positivamante o negativamente, se la successione delle somme n-esime diverge a \(+\infty\). Diverge negativamente, se la successione delle somme parziali \(-\infty\).






\end{document}
